\documentclass[11pt,a4paper]{article}

\usepackage[utf8]{inputenc}
\usepackage[T1]{fontenc}
\usepackage[french]{babel}
\usepackage{lmodern}
\usepackage{microtype}
\usepackage{geometry}
\usepackage{amsmath,amssymb,mathtools}
\usepackage{booktabs}
\usepackage{array}
\usepackage{hyperref}
\usepackage{xcolor}
\usepackage{enumitem}
\usepackage{tcolorbox}

\geometry{margin=2.4cm}
\hypersetup{colorlinks=true,linkcolor=blue!50!black,urlcolor=blue!50!black,citecolor=blue!50!black}
\setlength{\parindent}{0pt}
\setlength{\parskip}{0.65em}
\setlist[itemize]{topsep=0.25em,itemsep=0.2em}
\setlist[enumerate]{topsep=0.25em,itemsep=0.2em}

\newcommand{\dd}{\mathrm{d}}
\newcommand{\lp}{\ell_{\rm P}}
\newcommand{\rhoP}{\rho_{\rm P}}
\newcommand{\rhoDE}{\rho_{\rm DE}}
\newcommand{\rhoC}{\rho_{\rm c}}
\newcommand{\LH}{L_{\rm H}}
\newcommand{\EP}{E_{\rm P}}
\newcommand{\EH}{E_{\rm H}}
\newcommand{\Mpl}{M_{\rm P}}
\newcommand{\Order}{\mathcal{O}}
\newcommand{\rhoreg}{\rho_{\rm reg}}
\newcommand{\rhobh}{\rho_{\rm BH}}

\title{Relations fondamentales UV/IR, cellule GUP et energie noire\\
\large Verification mathematique et equations candidates}
\author{Note de travail}
\date{22 mai 2026}

\begin{document}
\maketitle

\begin{abstract}
Cette note rassemble des identites exactes et des equations candidates reliant
la cellule de phase deformee de type GUP, l'echelle de Planck, l'echelle de
Hubble et l'energie noire. Le resultat robuste est
\[
\rhoDE=\Omega_\Lambda\rhoC
=\frac{3\Omega_\Lambda}{8\pi}\rhoP\left(\frac{\lp}{\LH}\right)^2.
\]
Cette relation montre que le facteur $10^{-122}$ est un rapport d'aires
Planck--Hubble. Le regulateur GUP seul rend l'energie du vide finie, mais
d'ordre planckien si $\beta_0\sim1$. La piste mathematiquement coherente est
donc une double contrainte : regularisation UV par la cellule GUP et projection
IR par une borne gravitationnelle ou holographique.
\end{abstract}

\tableofcontents
\newpage

\section{Statut des formules}

On distingue trois niveaux.
\begin{enumerate}
  \item Les identites exactes, qui decoulent des definitions de $\rhoC$, $\LH$,
  $\rhoP$ et $\lp$.
  \item Les relations etablies dans des cadres connus : borne de trou noir,
  entropie de Bekenstein--Hawking, thermodynamique de Gibbons--Hawking et
  contrainte UV/IR de Cohen--Kaplan--Nelson.
  \item Les equations candidates, qui combinent le regulateur GUP avec une
  projection gravitationnelle globale.
\end{enumerate}

La note ne pretend donc pas prouver une theorie fondamentale. Elle isole les
relations qui restent exactes et celles qui peuvent servir de programme de
recherche.

\section{Echelles fondamentales}

Les grandeurs de Planck sont
\begin{align}
\lp &= \sqrt{\frac{\hbar G}{c^3}},&
t_P &= \sqrt{\frac{\hbar G}{c^5}},&
\EP &= \sqrt{\frac{\hbar c^5}{G}},\\
\Mpl &= \sqrt{\frac{\hbar c}{G}},&
\rhoP &= \frac{\EP}{\lp^3}=\frac{c^7}{\hbar G^2}.
\end{align}

L'echelle de Hubble et la densite critique d'energie sont
\begin{equation}
\LH=\frac{c}{H_0},
\qquad
\rhoC=\frac{3c^2H_0^2}{8\pi G}.
\end{equation}

La densite d'energie noire observee est
\begin{equation}
\rhoDE=\Omega_\Lambda\rhoC.
\end{equation}

\section{Identite Planck--Hubble--energie noire}

Comme $H_0=c/\LH$,
\begin{equation}
\rhoC=\frac{3c^4}{8\pi G\LH^2}.
\end{equation}
Or
\begin{equation}
\rhoP\lp^2
=\frac{c^7}{\hbar G^2}\frac{\hbar G}{c^3}
=\frac{c^4}{G}.
\end{equation}
Donc
\begin{equation}
\boxed{
\rhoDE
=\Omega_\Lambda\rhoC
=\frac{3\Omega_\Lambda}{8\pi}\rhoP
\left(\frac{\lp}{\LH}\right)^2.
}
\label{eq:master_uv_ir}
\end{equation}

\begin{tcolorbox}[colback=blue!3!white,colframe=blue!45!black,title=Lecture]
La petitesse de l'energie noire peut etre reecrite comme un rapport d'aires :
aire de Planck sur aire de Hubble. Cette phrase est une reformulation exacte,
pas encore une explication dynamique.
\end{tcolorbox}

\section{Borne de trou noir}

Une sphere de rayon $L$ et de densite d'energie $\rho$ contient
\begin{equation}
E=\rho\frac{4\pi}{3}L^3.
\end{equation}
L'energie d'un trou noir de rayon de Schwarzschild $L$ vaut
\begin{equation}
E_{\rm BH}=\frac{c^4L}{2G}.
\end{equation}
La condition de non-effondrement $E\leq E_{\rm BH}$ donne
\begin{equation}
\boxed{
\rho\leq\rhobh(L)=\frac{3c^4}{8\pi G L^2}.
}
\label{eq:rho_bh}
\end{equation}

Pour $L=\LH$, on retrouve exactement
\begin{equation}
\rhobh(\LH)=\rhoC,
\qquad
\rhoDE=\Omega_\Lambda\rhobh(\LH).
\end{equation}

\section{Echelle d'energie de l'energie noire}

Definissons
\begin{equation}
E_{\rm DE}=\left[\rhoDE(\hbar c)^3\right]^{1/4}.
\end{equation}
En utilisant \eqref{eq:master_uv_ir},
\begin{equation}
\boxed{
E_{\rm DE}
=
\left(\frac{3\Omega_\Lambda}{8\pi}\right)^{1/4}
\sqrt{\EP\EH},
\qquad
\EH=\hbar H_0.
}
\label{eq:E_de}
\end{equation}

L'echelle meV de l'energie noire est donc la moyenne geometrique entre
l'energie de Planck et l'energie de Hubble, a un facteur cosmologique pres.

\section{Regulateur GUP}

Dans le modele de cellule de phase deformee,
\begin{equation}
\dd N_{\rm GUP}
=
\frac{\dd^3x\,\dd^3p}{h^3(1+\beta p^2)^3}.
\end{equation}
Pour une energie de point zero massless,
\begin{equation}
\rhoreg
=
\frac{g}{h^3}4\pi\int_0^\infty
\frac{\frac12 pc\,p^2\dd p}{(1+\beta p^2)^3}.
\end{equation}
Comme
\begin{equation}
\int_0^\infty\frac{p^3\dd p}{(1+\beta p^2)^3}
=\frac{1}{4\beta^2},
\end{equation}
on obtient
\begin{equation}
\boxed{
\rhoreg
=
\frac{g\rhoP}{16\pi^2\beta_0^2},
\qquad
\beta=\beta_0\frac{\lp^2}{\hbar^2}.
}
\label{eq:rho_reg}
\end{equation}

Le GUP regularise l'UV. Mais pour $\beta_0\sim1$, la densite reste d'ordre
planckien.

\section{Projection holographique candidate}

La combinaison minimale des deux contraintes est
\begin{equation}
\boxed{
\rho_{\rm grav}(L)
=
\min\left[
\frac{g\rhoP}{16\pi^2\beta_0^2},
\frac{3c^4}{8\pi G L^2}
\right].
}
\label{eq:min_density}
\end{equation}

Cette equation est candidate : elle ne decoule pas du GUP seul. Elle formalise
l'idee que les modes UV peuvent etre regularises localement, tandis que la
gravite borne le nombre de degres de liberte independants dans une region de
taille $L$.

Le point de croisement des deux branches est obtenu par
\begin{equation}
\frac{g\rhoP}{16\pi^2\beta_0^2}
=
\frac{3c^4}{8\pi G L_\times^2}.
\end{equation}
En utilisant $\rhoP\lp^2=c^4/G$,
\begin{equation}
\boxed{
L_\times
=
\beta_0\sqrt{\frac{6\pi}{g}}\,\lp.
}
\label{eq:L_cross}
\end{equation}

Pour $\beta_0=1$ et $g=1$, $L_\times\simeq4.34\,\lp$. A toute echelle
macroscopique, la branche gravitationnelle domine donc largement.

\section{Facteur de projection}

Sur la branche gravitationnelle,
\begin{equation}
\rho_{\rm grav}(L)=\mathcal{F}(L)\rhoreg.
\end{equation}
Alors
\begin{equation}
\boxed{
\mathcal{F}(L)
=
\frac{6\pi\beta_0^2}{g}
\left(\frac{\lp}{L}\right)^2.
}
\label{eq:projection_factor}
\end{equation}

Le facteur minuscule vient donc du rapport d'aires $\lp^2/L^2$. Pour $L=\LH$,
il est de l'ordre de $10^{-122}$ si $\beta_0$ est d'ordre unite.

\section{\texorpdfstring{Lien avec la note sur $\beta_0^{(\Lambda)}$}{Lien avec la note sur beta0 Lambda}}

La note precedente imposait directement $\rhoreg=\rhoDE$. Cela donnait
\begin{equation}
\beta_0^{(\Lambda)}
=
\left(\frac{g}{6\pi\Omega_\Lambda}\right)^{1/2}
\frac{1}{H_0t_P}.
\end{equation}
L'energie GUP associee et l'echelle $E_{\rm DE}$ verifient alors
\begin{equation}
\boxed{
\frac{E_{\rm GUP}^{(\Lambda)}}{E_{\rm DE}}
=
\left(\frac{16\pi^2}{g}\right)^{1/4}
=
\frac{2\sqrt{\pi}}{g^{1/4}}.
}
\label{eq:energy_ratio}
\end{equation}

Pour $g=1$, le rapport vaut $3.5449$. Cette equation explique pourquoi
$E_{\rm GUP}^{(\Lambda)}\simeq7.94$ meV alors que
$E_{\rm DE}\simeq2.24$ meV.

\section{Thermodynamique d'horizon}

Pour un horizon de rayon $L$, la temperature de Gibbons--Hawking est
\begin{equation}
T_H=\frac{\hbar c}{2\pi k_B L}.
\end{equation}
L'entropie de Bekenstein--Hawking est
\begin{equation}
S_H=\frac{k_B A}{4\lp^2}
=\frac{k_B\pi L^2}{\lp^2}.
\end{equation}
Le produit $T_HS_H$ donne
\begin{equation}
T_HS_H
=
\frac{\hbar cL}{2\lp^2}
=
\frac{c^4L}{2G}.
\end{equation}
En divisant par $V=4\pi L^3/3$, on retrouve
\begin{equation}
\boxed{
\rho_H(L)=\frac{3c^4}{8\pi G L^2}.
}
\end{equation}

La meme densite apparait donc par deux voies : borne de trou noir et
thermodynamique d'horizon. Les signes et facteurs de premier principe dependent
du type d'horizon, mais l'echelle $L^{-2}$ est robuste.

\section{Holographie : volume GUP contre aire}

Le comptage volumique donne schematiquement
\begin{equation}
N_{\rm volume}(L)\sim\left(\frac{L}{\lp}\right)^3.
\end{equation}
La gravite impose plutot
\begin{equation}
\frac{S_{\rm BH}}{k_B}
=
\frac{A}{4\lp^2}
=
\frac{\pi L^2}{\lp^2}.
\end{equation}
Donc
\begin{equation}
\boxed{
\frac{S_{\rm BH}}{k_B}
\sim
N_{\rm volume}^{2/3}.
}
\end{equation}

Une formulation candidate est
\begin{equation}
\boxed{
N_{\rm grav}(\Omega_L)
=
\min\left[
N_{\rm GUP}(\Omega_L),
\frac{A(\partial\Omega_L)}{4\lp^2}
\right].
}
\end{equation}

\section{Relation de Weinberg cosmologique}

Une masse construite avec $\hbar,G,c,H_0$ est
\begin{equation}
\boxed{
m_W=\left(\frac{\hbar^2H_0}{Gc}\right)^{1/3}.
}
\label{eq:weinberg_mass}
\end{equation}
Elle equivaut, a des facteurs d'ordre unite pres, a
\begin{equation}
\lambda_C^2\sim r_g\LH,
\qquad
\lambda_C=\frac{\hbar}{mc},
\qquad
r_g=\frac{Gm}{c^2}.
\end{equation}
Cette relation est une echelle numerologique robuste, pas une prediction
dynamique du GUP. Elle signale un autre melange quantique--gravitation--Hubble.

\section{Acceleration cosmologique}

Avec
\begin{equation}
\Lambda=\frac{3\Omega_\Lambda H_0^2}{c^2},
\end{equation}
il faut distinguer
\begin{equation}
\boxed{
a_\Lambda=c^2\sqrt{\Lambda}
=cH_0\sqrt{3\Omega_\Lambda},
}
\end{equation}
et l'echelle de Sitter usuelle
\begin{equation}
\boxed{
a_{\rm dS}=c^2\sqrt{\frac{\Lambda}{3}}
=cH_0\sqrt{\Omega_\Lambda}.
}
\end{equation}
Pour les comparaisons de type MOND, la convention naturelle est souvent
\begin{equation}
\boxed{
a_0\sim\frac{a_{\rm dS}}{2\pi}
=
\frac{c^2}{2\pi}\sqrt{\frac{\Lambda}{3}}.
}
\end{equation}

\section{Equation cosmologique candidate}

Une ecriture effective inspiree par l'holographie est
\begin{equation}
\boxed{
H^2
=
\frac{8\pi G}{3c^2}\rho_m
+\eta(L)\frac{c^2}{L^2}.
}
\label{eq:friedmann_candidate}
\end{equation}
Elle correspond a une densite effective
\begin{equation}
\rho_{\rm eff}(L)
=
\eta(L)\frac{3c^4}{8\pi G L^2}.
\end{equation}

Si $L=c/H$, le terme est proportionnel a $H^2$ et peut se comporter comme une
renormalisation de $G$. Pour obtenir une energie noire avec $w\simeq-1$, il
faut une longueur IR dynamique ou non locale, par exemple l'horizon des
evenements futur dans les modeles d'energie noire holographique.

\section{Synthese}

Les relations les plus solides sont
\begin{align}
\rhoDE
&=
\frac{3\Omega_\Lambda}{8\pi}\rhoP
\left(\frac{\lp}{\LH}\right)^2,\\
E_{\rm DE}
&=
\left(\frac{3\Omega_\Lambda}{8\pi}\right)^{1/4}
\sqrt{\EP\EH},\\
\rhobh(L)
&=
\frac{3c^4}{8\pi G L^2},\\
\rho_{\rm grav}(L)
&=
\min\left[
\frac{g\rhoP}{16\pi^2\beta_0^2},
\frac{3c^4}{8\pi G L^2}
\right],\\
L_\times
&=
\beta_0\sqrt{\frac{6\pi}{g}}\,\lp,\\
m_W
&=
\left(\frac{\hbar^2H_0}{Gc}\right)^{1/3},\\
a_{\rm dS}
&=
c^2\sqrt{\frac{\Lambda}{3}}.
\end{align}

\begin{tcolorbox}[colback=yellow!5!white,colframe=yellow!40!black,title=Conclusion structurante]
\[
\text{GUP regularise l'UV}
\quad+\quad
\text{horizon projette l'IR}
\quad\Longrightarrow\quad
\text{residu holographique fini.}
\]
\end{tcolorbox}

Le probleme residuel n'est pas seulement le facteur $10^{-122}$. Il devient le
probleme plus precis de la dynamique d'horizon : pourquoi le bon $L$ est
selectionne, pourquoi le coefficient vaut aujourd'hui $\Omega_\Lambda$, et
pourquoi le tenseur effectif possede une equation d'etat proche de $w=-1$.

\begin{thebibliography}{9}
\bibitem{CODATA2022}
E. Tiesinga, P. J. Mohr, D. B. Newell and B. N. Taylor,
``CODATA recommended values of the fundamental physical constants: 2022'',
\emph{Reviews of Modern Physics} \textbf{97}, 025002 (2025).

\bibitem{Planck2018}
N. Aghanim et al. (Planck Collaboration),
``Planck 2018 results. VI. Cosmological parameters'',
\emph{Astronomy \& Astrophysics} \textbf{641}, A6 (2020),
arXiv:1807.06209.

\bibitem{CKN1999}
A. G. Cohen, D. B. Kaplan and A. E. Nelson,
``Effective Field Theory, Black Holes, and the Cosmological Constant'',
\emph{Physical Review Letters} \textbf{82}, 4971 (1999),
arXiv:hep-th/9803132.

\bibitem{Li2004}
M. Li,
``A Model of Holographic Dark Energy'',
\emph{Physics Letters B} \textbf{603}, 1--5 (2004),
arXiv:hep-th/0403127.

\bibitem{Bekenstein1973}
J. D. Bekenstein,
``Black Holes and Entropy'',
\emph{Physical Review D} \textbf{7}, 2333 (1973).

\bibitem{GibbonsHawking1977}
G. W. Gibbons and S. W. Hawking,
``Cosmological event horizons, thermodynamics, and particle creation'',
\emph{Physical Review D} \textbf{15}, 2738 (1977).

\bibitem{Kempf1995}
A. Kempf, G. Mangano and R. B. Mann,
``Hilbert space representation of the minimal length uncertainty relation'',
\emph{Physical Review D} \textbf{52}, 1108 (1995),
arXiv:hep-th/9412167.
\end{thebibliography}

\end{document}
